\documentclass[../main.tex]{subfiles}
\begin{document}
%==========================================TIÊU ĐỀ TẬP========================================
\begin{table}[h]
    \begin{adjustwidth}{-1cm}{}
        \begin{tabular}{>{\centering\arraybackslash}p{6cm}|>{\raggedright\arraybackslash}p{14cm}}
            \multirow{5}{6cm}
            % Cột bên trái: ÉPISODE và số tập
            {\\[-40pt]
                \flaregothic\fontsize{20pt}{20pt}\selectfont \centering
                \color{eptype}HISTOIRE\color{black}\\
                \gangofthree\fontsize{72pt}{72pt}\selectfont
                \color{epnum}1
                \\[18pt]
                % \flaregothic\fontsize{14pt}{14pt}\selectfont
                % \color{epattr}BIENVENUE, \uline{\color{Banpire} Mel} !
                % \flaregothic\fontsize{18pt}{18pt}\selectfont
                % \color{epattr}ÉPISODE PERDU
                \color{black}
            }
             &
            % Dòng thứ nhất: tên tập tiếng Nhật
            {\fotskip\fontsize{18pt}{18pt}\selectfont 「テト」って\ruby{何}{なに}？
            }                         \\[6pt]
            % Dòng thứ hai: tên tập tiếng Anh
             & {\cafeteria\fontsize{18pt}{18pt}\selectfont What is Tet?}                         \\[4pt]
            % Dòng thứ ba: tên tập tiếng Việt
             & {\vnmsans\fontsize{18pt}{18pt}\selectfont Tết là gì?}                             \\[8pt]
            % Dòng thứ tư: Nhân vật xuất hiện
             & {\brian{\fontsize{14pt}{14pt}\selectfont Track used: Lunar Zoo Year - Grasswalk}}
        \end{tabular}
    \end{adjustwidth}
\end{table}
%===========================================NỘI DUNG==========================================
\color{black}

\begin{center}
    (Tiếp tục từ {\flaregothic ÉPISODE 38 - SPÉCIAL {\ab V}}, phần {\abv Truyện phụ}, quyển số Ba.)
\end{center}

\pov{Hikawa Kuon}

Như tôi đã nói trước đó, hôm nay đã là 29 Tết. Ngày kia là sẽ bước sang năm mới Âm lịch ở Kawachi. Đáng lẽ cả gia đình tôi sẽ đón một cái Tết Nguyên Đán bình thường như bao cái Tết khác trong suốt hơn hai mươi nồi bánh chưng xanh của tôi.

Nhưng không, lần này lại khác. Và lý do thì thật ``không thể tin được'': toàn bộ \hll\ quyết định biến quê nhà của tôi thành phim trường để quay một bộ phim tài liệu ``trải nghiệm Tết \ruby{Etsunan}{Việt Nam}\ độc lạ. Không phải một vài người đâu, mà là cả công ty kéo quân đến nhà tôi. Tôi thề, nếu bạn là tôi, bạn sẽ hiểu cảm giác từ ngỡ ngàng chuyển sang bối rối rồi cuối cùng là\dots\ bất lực.

Làm sao tôi biết? Vì cách đây vài tiếng, ông giám đốc YAGOO, hay tôi vẫn gọi thân mật là Tanigo-san, đã gọi điện thông báo một cách rất nhẹ nhàng nhưng cũng vô cùng ``sát thương'': ``Hikawa-san, sẵn sàng tiếp đón tất cả idol nhà mình nhé. Họ đang trên đường rồi!''

``Tất cả idol ư? Nghĩa là khoảng hai mươi người ạ!?'' tôi hoảng hốt thốt lên, nhưng đáp lại chỉ là giọng cười đầy ý tứ.

``Ừ, tầm đó thôi. Mà đừng lo, họ mang theo quà nữa mà!''

Quà? Tôi không chắc món quà nào có thể bù đắp cho sự hỗn loạn sắp xảy ra đây. Đáng lẽ những ngày này tôi chỉ cần lo vài thứ đơn giản như sắm Tết, dọn nhà, phụ bố mẹ chuẩn bị mâm cỗ. Nhưng không, giờ đây lại có thêm nhiệm vụ tiếp đón ``đoàn khách quý'' bất đắc dĩ này. Còn bố mẹ tôi? Họ lại là những người hào hứng nhất, đặc biệt là bố tôi. ``\vie{Ui dồi, thế thì vui nhà vui cửa. Chứ mỗi năm chỉ có ba người thì buồn quá!}''

Hơi thở dài, tôi tặc lưỡi, Thôi thì cũng là cơ hội quảng bá văn hóa Kawachi, coi như đôi bên cùng có lợi vậy. Nhưng trong lòng, tôi vẫn không khỏi rùng mình nghĩ tới viễn cảnh hơn hai mươi nàng hề\dots\ à nhầm, nàng idol xuất hiện trước cửa nhà.

Và rồi, như để chứng minh dự cảm đó là thật, một âm thanh kỳ lạ phát ra từ phía phòng khách. Tôi quay lại và đứng hình mất ba giây. Cánh cửa thần kỳ xuất hiện ngay trên bức tường cạnh cửa sổ, đối diện chiếc vô tuyến.

Đúng vậy, một cánh cửa phát sáng, trông như bước ra từ truyện tranh D*r*e*on. Tôi chỉ vừa kịp há hốc miệng thì âm thanh ồn ào vang lên từ phía bên kia.

Tôi thề, họ không chỉ sẵn sàng mà còn có vẻ rất\dots\ ``nhiệt tình'' với sự kiện này. Mỗi người một túi đồ lỉnh kỉnh, tiếng cười nói vang cả căn nhà. Bố mẹ tôi thì vui như tết, còn tôi chỉ biết ôm đầu lẩm bẩm: ``\textit{Bắt đầu rồi đây\dots\ Cái Tết không giống ai năm nay của mình.}''

Tôi dám chắc họ đã chuẩn bị sẵn sàng, thậm chí là rất đầy đủ cho một trải nghiệm thuần Kawachi mới lạ này đây.

\il

\begin{center}
    \pov{-}
    \uline{Văn phòng\\Trước đó một tiếng.}
\end{center}

Trái ngược với suy nghĩ của cậu con trai, tất cả các thành viên vẫn đang tất bật chuẩn bị cho ``kỳ nghỉ'' (thực chất là buổi quay phim) Tết Nguyên Đán kéo dài bốn ngày.

Ý tưởng ban đầu là của Mio, người đã đề xuất quay một bộ phim tài liệu giới thiệu văn hóa Tết Etsunan, dựa theo những gì họ đã thấy trong thước phim ở phần credits. Nhưng việc chọn nhà Kuon làm địa điểm chính lại là quyết định bất ngờ của cô nàng Thần tượng Sora.

Đề xuất này không chỉ khiến các thành viên khác choáng váng, mà ngay cả A-chan, người luôn đồng hành cùng Sora trong các dự án táo bạo, cũng phải ngạc nhiên. ''Sora-san, không phải em thường rất thận trọng với những kế hoạch lớn sao? Sao lần này em lại quyết định gấp thế?'' bạn của cô hỏi, vẫn chưa hoàn toàn tin vào những gì mình vừa nghe.

Mio cũng không giấu được sự bất bình nhẹ nhàng, dù vẫn còn khó tin: ``Mà này, Sora-senpai, sao chị đưa ra tuyên bố gấp gáp thế!?''

``Bọn em đã chuẩn bị gì cho chuyến đi chơi này đâu!'' Fubuki thêm vào, vẻ mặt đầy lo lắng. A-chan chỉnh lời ngay lập tức: ``Chúng ta đi quay phim chứ không phải đi chơi đâu, Shirakami-san.''

``Nhưng mà Mio-sha nói cũng có lý đó. Chúng ta còn chưa biết văn hóa ăn Tết nhà anh ấy như thế nào mà!'' Subaru nói, giọng lộ vẻ bối rối.

``Ít nhất chúng ta cũng phải tìm hiểu trước chứ, ne!'' Miko tán thành, tay chống hông, ánh mắt đầy nghi hoặc hướng về phía Sora.

Nàng Đeo kính lắc đầu, khoanh tay lại: ``Cơ mà chị thấy rất ngạc nhiên là YAGOO-sama cũng đã đồng ý cho chuyện này rồi\dots\ Chẳng biết ông ấy nghĩ gì nữa.''

Nàng Thần tượng chỉ cười: ``Có gì đâu, ông ấy đã đồng ý thế rồi thì ta cứ triển khai thôi!''

Nàng Vịt nhíu mày: ``Nhưng vấn đề là ta sẽ chuẩn bị những gì và chuẩn bị như thế nào?'' Miko nhấn mạnh: ``Chúng ta mới chỉ biết là nhà anh ấy cũng ăn Tết như chúng ta thôi, ne. Nhưng Tết nhà anh ấy là gì thì chúng ta còn chưa biết\dots\'', đi qua đi lại trong phòng như thể đang cố tìm ra câu trả lời.

Một lúc sau, Matsuri ngồi bệt xuống ghế, hai tay chống cằm, bỗng hỏi một câu khiến tất cả mọi người ngớ ra: ``Vậy, Tết (Nguyên Đán) là gì?'' Câu hỏi của cô khiến cả phòng chìm vào im lặng. Từng ánh mắt dần quay sang nhau, nhưng chẳng ai nói được lời nào. Họ cũng không có câu trả lời rõ ràng.

Roboco phá tan bầu không khí tĩnh lặng bằng việc thử đoán về định nghĩa của nó: ``Chị nghĩ chắc là một vocaloid\footnote{Công nghệ tổng hợp giọng hát được ghi âm từ người thật (thường là các diễn viên hoặc ca sĩ).} nào đó chăng?''

``Kiểu. Kasane Teto\footnote{Trong tiếng Nhật, ``Tết'' (テト) đọc là ``teto''. Kasane Teto mà họ nhắc tới cũng là người hát bài \textit{Ochame Kinou} từ năm 2010, vốn được \hll\ cover lại vào 4/2020.} à?'' Shion hỏi lại, nheo mắt nghi ngờ, ``Em thấy hai cái này chẳng ăn khớp nhau lắm ấy.''

``Chị nghĩ là vậy\dots?'' nàng Người máy trả lời, vẻ mặt không mấy chắc chắn.

Nhắc đến vocaloid tóc đỏ hai bên xoắn ốc, Aqua và Matsuri bỗng đứng bật dậy đồng thanh hát bài \textit{Ochame Kinou}, bài hát mà cả hai vừa thu âm để mừng kỷ niệm 10 năm.

\begin{dialogue}
    \speak{Aqua} いつでも I love you, \ruby{君}{きみ}に Take kiss me{\fotskip \~{}} \\ (Lúc nào tớ cũng love you [yêu cậu], nên cậu hãy mau take [nhận lấy] rồi kiss me [hôn tớ đi]\~{})
    \speak{Matsuri} \ruby{忘}{わす}れられないから　\ruby{僕}{ぼく}の\ruby{大}{だい}\ruby{事}{じ}な\ruby{メ}{め}\ruby{モ}{も}\ruby{リ}{り}ー{\fotskip \~{}} \\ (Tớ không thể quên được những ký ức quý giá này đâu\~{})
\end{dialogue}

Mio, người vừa ngồi xuống để viết vài ghi chú, nghe thấy liền đổ mồ hôi hột: ``Tôi không nghĩ chỉ vì một người mà anh ấy sẵn sàng nghỉ gần một tháng đâu. Nó hẳn phải là một thứ gì đó to lớn hơn thế\dots''

Căn phòng lại rơi vào trạng thái im lặng đầy suy tư, chỉ trừ hai nàng idol vẫn đang say sưa hát.

Bất chợt, họ nhận được câu trả lời từ nàng Dơi, người đang chăm chú tra cứu trên máy tính. Mel dừng gõ phím, quay ghế lại và nói: ``Chị tra trên mạng thì nó giải thích là một lễ hội truyền thống ở vùng Etsunan, một đất nước nằm ở khu vực phía Tây Nam đất nước ta. Giống với phong tục ăn Tết Thiên Bảo ở Nhật trước năm 1873.''

``Sao lại là năm 1873?'' nàng Cổ động dừng hát, ánh mắt mở to đầy thắc mắc. A-chan đẩy gọng kính, từ tốn giải thích: ``Chị nhớ không nhầm thì từ thời Minh Trị thứ sáu người ta đã bỏ tục đó rồi. Nhật Bản chuyển sang dùng Dương lịch để hội nhập với phương Tây.''

Matsuri ngả người ra ghế, vẻ như đang tự trách bản thân. ``Ồ\dots\ Sao em lại không biết nhỉ?''

Nghe vậy, nàng Phù thủy lém lỉnh buông lời trêu chọc, môi nhếch lên nụ cười tinh nghịch: ``Dám chắc là chị ngủ gật giờ lịch sử luôn.'' Đáp lại lời khịa từ nàng tóc trắng là một nụ cười trừ từ Matsuri, người rõ ràng không thể chối cãi.

Trong khi đó, những người khác bắt đầu dồn sự chú ý vào màn hình của Mel, cố tìm hiểu thêm về lễ hội mới mẻ này. ``Cụ thể `Tết' bên đó là như thế nào ạ?'' Subaru nghiêng người về phía cô nàng Dơi, cố nhìn rõ hơn những gì đang hiển thị.

``Trong này cũng ghi rằng nó được tổ chức khá giống với năm mới bên ta. Họ cũng có đi đền chùa, cũng nhân dịp này để thăm gia đình họ hàng, và có những lễ hội gắn liền với những ngày này nữa,'' nàng Dơi giải thích, tay chỉ vào vài dòng thông tin trên màn hình.

``Những ngày này? Ý chị là sao?'' Aqua nhíu mày, giọng đầy nghi hoặc.

Mel vừa nói vừa lướt mắt đọc thêm các ghi chú: ``Họ ăn Tết dài lắm, thường từ 29 tháng Chạp đến hết mùng Ba tháng Giêng. Nhưng thực tế họ nghỉ từ khoảng 23 tháng Chạp đến hết Rằm tháng Giêng cơ.''

``Ể? `Chặp'? `Riêng'? Tháng gì nghe lạ vậy, peko?'' Pekora gãi đầu, đôi tai thỏ lắc lư theo nhịp câu hỏi của cô.

Mel trả lời, kiên nhẫm giải thích thêm: ``Tháng 12 với tháng 1 Âm lịch đó. Nước ta cũng đã từng sử dụng lịch Âm như họ rồi.''

``Vậy cụ thể là bao nhiêu ngày vậy?'' Aqua lại hỏi. Ayame hớn hở đáp, như thể cô vẫn còn năng lượng sau cuộc họp về nội dung bộ phim từ trước đó: ``Tôi nhớ đâu đó hơn hai mươi ngày á.''

``Họ tổ chức lễ hội đó dài đến thế ư!? Ngày Tết chúng ta chỉ được nghỉ có ba ngày thôi á!'' nàng Cổ động không giấu nổi sự ngạc nhiên.

``Cũng tùy từng vùng thôi, dazo'', nàng Oni nhún vai, như thể mọi chuyện đều đơn giản với cô ấy vậy. Aqua và Matsuri nghe vậy, càng tỏ vẻ khó hiểu hơn.

Fubuki chen ngang, mắt nhìn Mel đầy hiếu kỳ: ``Mà, Mel-chan, `Rằm tháng Giêng' là ngày gì vậy?'' Mel trả lời, giọng nhẹ nhàng như đang đọc một bài thơ: ``Ngày 15 tháng Giêng Âm lịch, ngày trăng tròn nhất trong tháng.''

``Thậm chí bên đó người ta còn nói vui là: `Còn Mùng là còn Tết mà' dazo,'' Ayame thêm vào, khóe môi cong nhẹ.

Cả nhóm nghe xong đều đờ người ra vài giây. Nàng Vịt gãi đầu, rồi lẩm bẩm: ``Còn Mùng là còn Tết? Ý là họ ăn Tết cả tháng à!?''

Pekora chắp tay sau đầu, vẻ mặt không giấu được sự ngạc nhiên, xen chút một chút ngưỡng mộ: ``Họ đúng là khỏe thật đấy, peko\dots\ Ăn chơi đến thế thì ai chơi lại được chứ?''

Matsuri, đôi mắt sáng rỡ, reo lên: ``Này, thế là mình có thể mượn cái `Rằm tháng Giêng' đó làm cớ để kéo dài chuyến đi này đúng không!?'' Mio nghe vậy, thở dài: ``Chúng ta đến đó để quay phim, chứ không phải đi nghỉ dưỡng đâu, Matsuri.''

Aqua cười khì, khi kế hoạch của hai người họ bị bại lộ: ``Nhưng mà nghe cũng hấp dẫn ghê.''

Subaru tiến tới người bạn Quỷ sừng đang ngáp ấy: ``Cơ mà, sao bà biết được mấy cái này vậy?'' Nàng Oni hãnh diện đáp: ``Tôi từng đi chơi Tết ở đó một lần rồi. Điều mà tôi thấy bất ngờ là mỗi vùng họ tổ chức khác lắm đó!''

``Em nói mỗi vùng một khác là thế nào?'' nàng Cáo trắng ngạc nhiên. Ayame nghiêng đầu, tay xoa cằm như đang cố nhớ lại: ``Đặc trưng Tết của hai miền vùng đó khác hẳn nhau á. Đây nè, miền Bắc thì Tết của họ thường rất lạnh, giống như mùa đông bên chúng ta. Còn miền Nam thì ấm áp, thậm chí nhiều năm còn nắng gắt ấy!''

Rushia ngơ ngác: ``Tại sao lại có sự khác biệt ngược đời này nhỉ?'' A-chan nhẹ nhàng giải thích, giọng như thể cô đang ở trong một buổi thuyết trình: ``Chị nghĩ là do sự khác biệt về địa lý ở vùng đất đó thôi. Miền Bắc đất nước họ theo chị thấy là nằm gần khu vực ôn đới nên có bốn mùa rõ rệt, còn miền Nam gần xích đạo hơn nên chỉ có hai mùa: mùa khô và mùa mưa.''

Nàng Sói đen gật gù: ``Đúng rồi! Em nhớ là từng đọc đâu đó rằng khí hậu của vùng này chia làm hai phần rõ rệt như vậy. Nhìn như em nhớ là nó được ngăn cách bởi dãy núi nào đó\dots''

Ayame tiếp lời: ``Vì thế, văn hóa ăn Tết của hai vùng đó cũng khác nhau lắm. Ví dụ nhé, miền Bắc thì có cây đào hồng, còn miền Nam lại chuộng cây mai vàng. Cả hai miền thì nhà nào cũng hay trưng cây quất, giống như cây tùng của bên mình. Rồi còn cây nêu - trông nó như cây tre trang trí thêm đồ vật để xua đuổi tà ma.''

``Dám chắc là họ cũng sẽ xua đuổi cả bà nữa nhỉ, Ayame?'' Subaru nhoẻn miệng cười.

``Dĩ nhiên là không rồi!'' nàng Oni phồng má lên. ``Tôi là Quỷ tốt mà!''

Nàng Dơi xoay chiếc ghế lại về phía họ, cười tủm tỉm: ``Nghe thú vị ghê! Mà bên đó cũng dùng màu đỏ làm chủ đạo để cầu may giống như chúng ta, đúng không?''

Nàng Quỷ sừng gật đầu, hồ hởi quay về chủ đề chính: ``Chính xác luôn! Rồi còn có các hoạt động truyền thống nữa.''

Nàng Pháp sư tò mò hỏi: ``Vậy họ có những hoạt động gì?''

``Ví dụ, đêm Giao thừa thì nhà nhà quây quần ăn bữa cơm tất niên. Em từng đến dự một bữa tất niên rồi, không khí ấm cúng lắm! Nói chung là còn nhiều thứ để tận hưởng nữa!'' Ayame nhanh nhảu đáp.

Nàng Vu nữ chớp mắt, tò mò hỏi: ``Vậy họ có đi đền chùa hay đi cầu Thần không, ne?''

``Có chứ'', nàng Oni đáp nhanh, ``người ta hay đi bốc quẻ đầu năm để lấy may, hoặc cầu nguyện cho sức khỏe và hạnh phúc. À, còn một cái hay nữa là họ đi xin chữ!''

Matsuri nghe đến đây thì mắt sáng rỡ: ``Xin chữ? Có phải là giống những buổi fan meeting xin chữ ký từ idol không?''

Trong đầu nàng Lễ hội bắt đầu hình dung cảnh người người xếp hàng dài, chen lấn để giành lấy chữ ký của chính mình. Cô cười mãn nguyện, tưởng tượng mình là tâm điểm của đám đông hâm mộ.

Mio vội gạt đi ý nghĩ đó: ``Em không nghĩ người ta dự đông đến mức độ đó đâu. Với cả em không nghĩ ý `xin chữ' là như thế\dots'' Nhưng dường như cô nàng Cổ động đó không nghe thấy. Cô vẫn mơ màng, vẻ mặt ngập tràn tự hào.

Subaru thở dài, liếc mắt nhìn Matsuri: ``Chị ấy thực sự muốn có danh vọng đến thế ư\dots?'' Miko thì che miệng cười khúc khích: ``Matsuri thật là\dots\ nhưng mà, xin chữ đúng là một phong tục độc đáo nhỉ? Em muốn thử quá!''

Nàng Cáo trắng xen vào, mắt sáng lên: ``Nhưng mà, chữ họ xin là gì vậy? Có giống kiểu bùa may mắn không?''

Ayame khẽ lắc đầu: ``Không hẳn đâu. Chữ thường là do các `ông đồ' - kiểu như các thầy đồ xưa bên mình - viết. Họ viết những chữ mang ý nghĩa tốt đẹp, như `Phúc' (\ruby{福}{ふく}), `Lộc' (\ruby{禄}{ろく}), hoặc `Thọ' (\ruby{寿}{しゅ}) chẳng hạn.''

Nàng Sói đen mỉm cười: ``Nghe như vậy, Tết của họ vừa có nét gần gũi với chúng ta, lại vừa độc đáo nữa. Nếu có cơ hội, chị cũng muốn thử trải nghiệm một lần.''

Nàng Vu nữ chắp tay, vẻ mặt quyết tâm: ``Vậy thì năm nay tôi phải tận dụng cơ hội này thật tốt! Ai muốn xin chữ của tôi thì cứ đăng ký trước ne!''

Tất cả mọi người đều bật cười với sự nhiệt tình của nàng tóc hồng kia. Nàng Quỷ sừng tiếp tục hào hứng kể: ``Trong lễ hội đó còn có nhiều thú vui lắm nha, đặc biệt là múa lân!''

Nàng Lễ hội nghiêng đầu khó hiểu, cô ghé tai Haato, nói nhỏ: ``\hl{Nè, lân là gì vậy?}''

Cô nàng từ đội cổ động bỗng tưởng tượng đến cảnh mình nhảy múa cùng một con sư tử thật, bộ lông xù xì của nó và đôi mắt rực lửa nhìn cô chằm chằm. Đáng sợ hơn, cô hình dung nó há miệng ngoạm lấy mình. Toàn thân của nàng Cổ động run lên bần bật: ``N-Người ta dũng cảm thật đấy\dots''

Subaru và Mio nhìn cô với những cặp mắt đầy khó hiểu. Nàng Sói đen thở dài, lấy tay bóp nhẹ sống mũi rồi lắc đầu: ``Ý nghĩa của `múa lân' không phải như thế đâu\dots\ Mà chị ấy đang nghĩ gì không biết nữa\dots''

Haachama, vẫn chưa hết tò mò, hỏi Ayame tiếp: ``Vậy còn phong tục thì sao?''

Ojou-san phấn khởi trả lời: ``Ồ, phong tục bên đó thú vị lắm! Vào 23 tháng Chạp, họ sẽ làm lễ tiễn Ông Công Ông Táo. Sau đó, mùng Một thì có tục xông đất, rồi còn cúng tổ tiên trên bàn thờ nữa!''

Fubuki ngạc nhiên: ``Ông Công\dots\ Ông Táo? Nghe lạ quá!'' Matsuri thì nhăn mặt: ``Xông\dots\ đất? Cúng thờ? Là sao cơ?'' Trong khi đó, Miko thì nghiêng đầu khó hiểu, mặc dù bản thân cô cũng đang làm ở trong đền.

Noel, với sự tưởng tượng phong phú của mình, bắt đầu nghĩ về viễn cảnh cô xông thẳng vào một cục đất đá to lớn. Chẳng mấy chốc, cô đã dùng sức mạnh để đập vỡ nó thành từng mảnh. ``\dots Có phải kiểu như vậy không?''

Nàng Thỏ nhìn cô, bật cười: ``Bà đang nghĩ linh tinh cái gì vậy, peko? Ai cũng xông đất kiểu đó thì rừng núi tan hoang hết à!''

Mel tiếp tục đọc những thông tin trên máy tính, rồi từ tốn giải thích: ``Chị nghĩ `xông đất' là việc người đầu tiên bước vào nhà ai đó trong ngày đầu năm mới để đem lại may mắn thôi, đúng không?''

Nàng Quỷ sừng gật đầu: ``Chắc vậy. Rồi sau đó người ta chuẩn bị bàn thờ cúng tổ tiên vào ngày Tết, để mời những người đã khuất về chung vui với gia đình.''

Miko chắp tay, giọng đầy thành kính: ``Kami-sama phù hộ cho họ mà, ne!'' Rushia đổ mồ hôi hột: ``Em không nghĩ vậy đâu\dots''

Mel hơi nghiêng đầu, vẻ băn khoăn: ``Nghe mọi người nói thì thấy phong tục này cũng sâu sắc ghê. Nhưng mà em nghĩ cơ bản nó cũng giống bên ta thôi. Cả hai bên đều có bữa cơm tất niên, đều đi đền chùa cầu may\dots''

Sora cười hiền, nhẹ nhàng hỏi: ``Nhắc đến bữa cơm thì bên đó thường có những món truyền thống gì, Ayame-chan?'' Ngay vừa khi nhắc đến đồ ăn, Noel gần như giật thót, chạy đến nàng Quỷ sừng, mắt sáng rực lên: ``Người ta có ăn thịt như em không!?''

Nàng Hiệp sĩ, mỗi khi nghe đến đồ ăn, luôn trở nên phấn khích như vớ được kho báu. Cô không ngừng lặp lại: ``Cho em ăn thử đi!'' khiến Choco và Shion phải vội vàng giữ lại để tránh cô vùng lên.

\begin{dialogue}
    \speak{Choco} *vỗ vai, cười gượng* Bình tĩnh lại nào, Noel-sama!
    \speak{Shion} Chị nghĩ người ta cũng sẽ ăn thịt như chúng ta thôi!
\end{dialogue}

Nhìn cô nàng phản ứng như trẻ con (lần nào cô cũng vậy), Pekora khoanh tay, thở dài: ``Lại bắt đầu rồi đấy, peko\dots''

Dù thế, cô nàng vẫn không hề nản lòng, tiếp tục giục: ``Nói đi, nói đi, nói đi ($ \times\ 2\pi $), Ayame-senpai! Người ta thường ăn gì? Có thịt gà không? Có thịt bò không? Có-''

Nàng Y tá siết nhẹ vai Noel, thì thầm: ``Noel-sama, kiềm chế lại cảm xúc của mình đi!'' Thấy hậu bối của mình đang nhìn cô với sự thèm thuồng, Ayame cười trừ, tiếp tục giải thích: ``Bên đó họ ăn nhiều loại thịt lắm, nhưng đâu chỉ có mỗi thế!''

Nàng Thỏ gật đầu đồng tình với cô: ``Ắt hẳn người ta cũng ăn rau củ chứ nhỉ, peko?''

Haato chắc nịch: ``Chị đám chắc luôn.'' Nàng Quỷ sừng thốt lên: ``Đương nhiên!'' rồi bắt đầu kể chi tiết: ``Các món truyền thống của hai miền cũng khác hẳn nhau! Miền Nam có thịt kho với trứng vịt (thịt kho hột vịt), bánh tét và củ kiệu. Miền Bắc thì có thịt lợn đông, bánh chưng, bánh giầy\footnote{Ở một số nơi, người ta gọi là \uline{bánh dầy}.} và dưa hành. Ngoài ra, họ còn ăn cá, thịt bò, và nhiều món khác nữa, tùy từng nhà!''

Okayu vuốt cằm, nhận xét: ``Cơ bản họ cũng ăn món mặn như ta nhỉ\~{}'' Korone hào hứng tiếp lời: ``Đúng vậy ha, Okayu\~{}''

Rushia nghiêng đầu với những cái tên mà cô chưa từng nghe bao giờ: ``Bánh chưng, bánh tét\dots\ Rồi bánh giầy\dots\ là thế nào vậy?''

Nàng Phù thủy phỏng đoán: ``Chị nghĩ chắc nó phải gắn liền với cái gì đó. Kiểu như, nó được tạo thành từ phép thuật gì đó mà ta không biết\dots''

Nàng Y tá bật cười: ``Cô không nghĩ là nó được tạo từ mấy thứ như vậy đâu, Shion-sama.''

``Vậy thì nó có thể là cái gì nhỉ?'', Matsuri tỏ vẻ thắc mắc, ``Có ai ở đây biết không?''

Không một ai biết câu trả lời. Các khái niệm thuần Etsunan này hẳn là một thứ gì đó khá lạ lẫm đối với họ.

Sau một hồi ngẫm nghĩ, bỗng Sora thốt ra câu trả lời: ``Hình như chị đã từng đọc về nó rồi thì phải\dots'' Người bạn của cô - A-chan - mỉm cười khích lệ: ``Em nói thử cho mọi người nghe xem nào.''

Cô chậm rãi giải thich, theo những gì mà cô biết: ``Em không nhớ chắc chắn, nhưng hình như có một người đã dâng lên vua hai món bánh chưng và bánh giầy đúng vào dịp đầu xuân. Trong đó, bánh chưng tượng trưng cho đất, còn bánh giầy tượng trưng cho trời. Nhà vua đồng ý và quyết định lấy hai món đó làm món ăn đặc trưng cho ngày Tết.''

Nàng Người máy nghiêng đầu, tròn mắt: ``Đất vuông? Trời tròn\dots? Ơ là sao??'' trên đầu cô đầy rẫy những dấu hỏi. Matsuri cũng bày tỏ đồng cảm với cô: ``Em cũng thấy khó hiểu như chị á, Roboco-san\dots''

Choco suy tư: ``Cô nghĩ quan niệm của người xưa là như vậy đấy. Họ hình dung đất vuông và trời tròn để nói lên sự hài hòa của vũ trụ.''

Subaru bật cười: ``Cũng giống như huyền thoại thỏ trên Mặt Trăng thôi à.''

Bỗng nhiên, mắt của Fubuki sáng lên như thể cô vừa nảy ra một ý tưởng: ``Chúng ta đã làm một tập phim nào về truyền thuyết hay cổ tích chưa nhỉ\dots?''

Mio ôm trán, ngán ngẩm: ``Fubuki, đừng có cho mọi người thêm ý tưởng điên rồ nữa đi. Bà thấy chúng ta còn chưa đủ loạn hay sao? Với cả ta đã chốt là làm về Tết rồi.''

Dường như không chỉ Rushia, nàng Cổ động và nàng Người máy tỏ vẻ khó hiểu. Noel cũng vậy: ``Nhưng nếu thế thì bánh tét là như thế nào?''

Mel tiếp tục ngồi tra cứu, mắt cô dán vào màn hình từ nãy đến giờ để tìm hiểu thêm về các phong tục Tết. Cô nhẹ nhàng lên tiếng: ``Chị tra trên mạng thấy nói rằng bánh tét là bánh chưng, nhưng được làm khuôn tròn, ở trong miền Nam mới có.''

Nàng Hiệp sĩ chớp mắt, tỏ vẻ ngờ vực: ``Ngoài cái đó ra thì bánh chưng với bánh tét khác gì nhau?''

Nàng Thỏ khoanh tay, thở dài: ``Chúng giống như nhau thôi, chỉ khác một cái là khuôn vuông, một cái khuôn tròn thôi, peko. Thông não một chút đi chứ.''

Noel gật gù, không lấy làm tức giận với câu cà khịa của cô. Cô hỏi ngay: ``Vậy nó làm từ gì?''

Ayame mỉm cười, lấy ngón tay ra đếm nguyên liệu: ``Chị nhớ hình như là có đậu xanh, lá dong, gạo nếp, thịt lợn-'' Nhắc đến thịt, nàng Hiệp sĩ lập tức không kìm được, phấn khích ngắt lời: ``\textbf{Có cả thịt nữa sao!?}''

Cô định lao lên và vô lấy nàng Quỷ sừng xuống để hỏi đầy kỹ càng, nhưng Choco và Shion đã nhanh tay giữ lại. Nàng Y tá kéo cô ngồi xuống, còn nàng Phù thủy giơ tay giả bộ tung phép để đe dọa.

\begin{dialogue}
    \speak{Choco} *thở dài* Đó, lại nữa rồi\dots\ Để người ta nói hết đã nào.
    \speak{Shion} Có cần tôi cho cô một phép bùa ngủ không đây hả, Noel!?
\end{dialogue}

A-chan nhìn hai cô nàng đang cố gắng dìm nàng Hiệp sĩ lực điền đó liền đổ mồ hôi hột, rồi quay sang Ayame cười trừ: ``Ờm\dots\ Nakiri-san, em nói tiếp đi.''

Nàng Quỷ sừng tiếp tục: ``Người ta làm nó từ đậu xanh, thịt lợn với gạo nếp. Sau đó bọc bên ngoài một lớp lá dong - em nghĩ tên nó là vậy - thắt nút lại, rồi đem đi luộc.'

Nàng Song tính gật gù, nhận xét: ``Nghe có vẻ đơn giản nhỉ.'' Nàng Oni bật cười: ``Và đây mới là phần hay nhất này. Canh nồi bánh chưng!''

Korone nghiêng đầu: ``Tại sao lại như vậy?'' Bạn của cô tò mò, kéo dài giọng: ``Cụ thể là thế nào vậy\~{}?''

Ayame hào hứng giải thích: ``Trong lúc chờ bánh được luộc chín, người ta hay làm nhiều trò để giết thời gian. Nào là đánh bài, nhảy múa xung quanh nồi bánh, hoặc kể chuyện vui. Nói chung là nhiều hoạt động lắm!''

Cặp Chó-Mèo bắt đầu hình dung ra cảnh tượng họ đang thực hiện những hoạt động quanh mội cái nồi được đun trên bếp lửa. Nhưng, trí tưởng tượng phong phú của cặp đôi này nhiều lúc lại trở nên hơi\dots\ quá đà.

\begin{dialogue}
    \speak{Korone} Ta có thể nhảy cà tưng quanh nó, đúng không ha, Okayu?
    \speak{Okayu} Thậm chí ta còn có thể uống trà với nó nữa\~{}
    \speak{Korone} Chơi đuổi bắt\dots
    \speak{Okayu} \dots hoặc thậm chí là tâm sự với nó ha\~{}
\end{dialogue}

Subaru nhìn hai người họ đang mơ mộng, thậm chí còn nhảy cà tưng quanh một chiếc nồi bánh chưng (tưởng tượng), liền đổ mồ hôi hột, lắc đầu: ``Chẳng biết họ đang nghĩ cái gì nữa\dots''

Ayame nghiêng đầu sang phải, nhắm mắt lại, rồi cố nhớ lại theo trải nghiệm mà cô biết: ``Tôi nghĩ chắc phải tầm 12 tiếng gì đó. Từ chiều tối hôm trước đến sáng sớm hôm sau.''

Nàng Dị giới trợn mắt ngạc nhiên: ``Sao lâu vậy?'' Nàng Tomboy cũng tỏ vẻ khó hiểu: ``Đến cả tôi cũng thấy hơi bất ngờ đấy.''

Nàng Đeo kính đăm chiêu, thử đưa ra một lý luận: ``Chị nghĩ chắc là để nó nhừ hẳn. Với cả, chị đoán người ta không chỉ luộc một cái mỗi mẻ đâu.''

Sora gật đầu đồng tình: ``Em dám chắc họ sẽ làm cả chục cái mỗi lần luộc như thế.''

Mio khoanh tay, mỉm cười: ``Nhưng nếu người ta không đem bán thì không chắc họ sẽ làm nhiều đến vậy đâu\dots''

Sau khi bàn tán một hồi về bánh chưng và canh nồi bánh, không khí vẫn rất sôi nổi. Haachama, luôn đầy năng lượng, liệt kê nhanh những món ăn mà cả hai bên thường có trong dịp năm mới:`` Món mặn này, món chay này\dots\ có chính, có phụ rồi\dots''

Cô quay sang hỏi nàng Ma cà rồng và Ojou-san: ``Thế còn tráng miệng và đồ ăn vặt ngày Tết bên đó có gì vậy?'' Ayame mỉm cười, vừa nói vừa đưa tay minh họa: ``Có đủ loại bánh kẹo, đều rất đặc biệt! Nhưng mứt là món ăn vặt truyền thống ngày Tết của họ bên đó! Kiểu như bánh mochi ở bên mình ấy.''

Cáo trắng ngẫm nghĩ, nhưng vẫn chưa hình dung được, liền hỏi: ``Nó trông như thế nào vậy?'' Cô đáp, tỏ vẻ hào hứng: ``Có nhiều loại lắm! Mứt dừa, mứt cà rốt, mứt lạc\dots''

Pekora nhướng mày, ngạc nhiên: ``Có cả cà rốt à, peko? Cơ mà, mút là gì vậy?'' Mel cười nhẹ, góp lời: ``Chị nghĩ chắc chúng giống như kẹo sấy khô thôi nhỉ.''

Rushia khoanh tay, cười nhẹ: ``Chắc có khi chúng ta ở đó hết Tết mà vẫn chưa trải nghiệm hết mấy món ăn bên đó quá nha.''

Ayame chọc ghẹo, nhìn nàng Hiệp sĩ: ``Cá luôn là Noel-chan sẽ thích cho mà xem.''

Sau khi được nghe tất cả những gì mà nàng Dơi, nàng Quỷ sừng và nàng Thần tượng tìm hiểu được về ngày lễ truyền thống của xứ Etsunan, mọi người giờ đã hiểu được phần nào về văn hóa của lễ hội này. Shion khoanh tay, khẽ gật đầu: ``Xem chừng vẫn còn nhiều thứ ta phải tìm hiểu về Tết nhỉ.''

Tuy thế, Ayame dường như chưa muốn dừng lại. Cô hào hứng nói tiếp: ``Ngày Tết bên đó, họ còn làm nhiều thứ khác nữa! Ví dụ như\dots''

Nàng Sói đen giơ tay ra hiệu, cắt ngang: ``Rồi rồi, chị nghĩ chúng ta cũng hiểu được sơ bộ về 'Tết' rồi. Khi nào sang nhà của Kuon-senpai thì ta có thể tìm hiểu kỹ hơn, được chứ?''

Nàng Đeo kính tán thành: ``Chị cũng nghĩ là nên trải nghiệm thực tế vẫn tốt nhất.'', để rồi nàng Thần tượng mỉm cười, gật đầu: Chứ ngồi kể thế này có khi chúng ta vẫn chưa hình dung được đâu.

Tất cả mọi người đều gật đầu đồng ý với suy nghĩ của nàng Thần tượng. Nàng Quỷ sừng buột miệng nhận xét: ``Nhưng mà nghe mọi người bàn tán sôi nổi về nó như thế chắc là cũng háo hức lắm nhỉ?''

\begin{dialogue}
    \speak{Fubuki} Đúng, đúng!
    \speak{Okayu} Bọn em muốn thử chơi ``Tết'' ở đó.
    \speak{Korone} Em cũng thế!
    \speak{Roboco} Ca mình mữa!
    \speak{Noel} Em muốn ăn thử đồ ăn bên đó!
    \speak{Matsuri} Em muốn được tới đó chơi!
    \speak{Mio} Chúng ta tới đó để quay phim chứ không phải đi nghỉ đâu.
    \speak{Fubuki} Dù gì chúng ta cũng sẽ cắm trại ở gần nhà \emph{anh ấy} mừ\~{} Nên cơ bản là cũng không khác gì nhau đâu.
    \speak{Mio} \dots Bà nói cũng phải.
\end{dialogue}

A-chan nhìn quanh, mỉm cười: ``Có vẻ như không ai phản đối gì nhỉ.'' Subaru nhún vai, cười nhẹ: ``Chúng ta đã biết sơ bộ về nó rồi, giờ chắc sẽ không thấy quá sốc nữa ta?''

Sora mỉm cười, ánh mắt lấp lánh: ``Cũng là dịp để biết ngày xưa chúng ta đã từng ăn Tết thế nào nữa!'' Ayame cười tươi, nhấn mạnh: ``Nhưng em nghĩ chúng ta sẽ còn đón chờ nhiều điều bất ngờ nữa đó. Chuẩn bị tinh thần đi nha\~{}''

Nàng Thần tượng giơ nắm đấm lên cao, cổ vũ cho mọi người: ``Nào, mọi người chuẩn bị trải nghiệm văn hóa Tết chưa?''

Tất cả đều hưởng ứng nhiệt tình, đồng thanh hô to: ``\textbf{Rồi! (お\~{}!!)}''

``\textbf{Vậy còn chờ gì nữa? Nhà của Hikawa-san, thẳng tiến!}''

``\textbf{Tiến!}''

A-chan bất chợt nhắc nhở: ``À mà, đừng quên chúng ta đến đó để quay phim nha!''

Tất cả cùng đáp, pha chút lém lỉnh: ``Rõ\~{}''

Toàn bộ gần hai mươi người từ \hll\ tập hợp tại văn phòng, nơi vốn đã quen thuộc nhưng lại chứa đựng những bí mật mà không phải ai cũng biết.

Hôm nay, các cô gái sẽ khám phá một trong những bí mật ấy.

Yuujin A, một trong những thành viên cấp cao nhất trong \textit{holoStaff}, dẫn đầu đoàn cùng với Tokino Sora. Trong tay cô là chiếc chìa khóa sơ mà Kuon từng đưa cho (phòng trường hợp khẩn cấp). Cô mở một cánh cửa dẫn vào phòng thu âm ngay bên trong văn phòng, nơi ít ai ngờ rằng lại có lối đi bí mật dẫn tới một điều bất ngờ.

Bên trong phòng thu, mọi thứ trông bình thường như mọi khi: micro, tai nghe, và một vài kịch bản còn dang dở trên bàn. Nhưng ở góc cuối căn phòng, đằng sau một chiếc tủ đựng thiết bị âm thanh, A-chan ấn nhẹ vào một công tắc ẩn. Chiếc tủ từ từ trượt sang một bên, để lộ một cánh cửa màu đỏ, trông rất khác biệt so với không gian xung quanh.

Tuyệt nhiên không có một bóng người, cả căn phòng trở nên yên tĩnh đến lạ thường. Cánh cửa đỏ, được bảo vệ cẩn thận, chính là lối đi giữa hai thế giới mà Kuon sử dụng để trở về Kawachi.

Mọi người xúm lại gần, ánh mắt tràn ngập sự hồi hộp và tò mò. Đứng trước cánh cửa kỳ diệu, nàng Thần tượng khẽ hít một hơi sâu, giọng dõng dạc như chuẩn bị bước vào một hành trình vĩ đại: ``Mọi người\dots\ Cùng đến vùng đất của anh quản lý chúng ta nào.''

Tất cả hai mươi cô gái đều gật đầu đồng thanh: ``\textbf{Ưm!}'' ``\textbf{Đi thôi!}''

Cô từ từ xoay nắm cửa và mở ra. Tất cả đều nín thở, hồi hộp nhưng không kém phần phấn khích. Khi cánh cửa mở hẳn ra, một luồng ánh sáng ấm áp ùa vào, như chào đón họ đến với thế giới mới.

Không kiềm chế được sự háo hức, các cô gái \hll\ nhanh chóng bước qua cánh cửa, để lại phía sau những tò mò dang dở và tiến vào một hành trình không ai ngờ tới\dots

\end{document}